\documentclass[../../../include/open-logic-section]{subfiles}

\begin{document}

\olfileid{sth}{cardinals}{hp}
\olsection{परिशिष्ट: ह्यूमचे तत्त्व}

\olref[cp]{sec} मध्ये आपण कँटरचे तत्त्व मांडले. ते असे:
\begin{align*}
  \card{A} = \card{B} & \text{ तेव्हा आणि केवळ तेव्हाच } A \approx B.
\intertext{आज ज्याला \emph{ह्यूमचे तत्त्व} म्हणतात त्याच्याशी हे खूप मिळतेजुळते आहे. ह्यूमचे तत्त्व असे:}
  \fregenum{x} {F(x)} = \fregenum{x}{G(x)} & \text{ तेव्हा आणि केवळ तेव्हाच } F \sim G
\end{align*}
येथे ‘$F \sim G$’ याचा अर्थ $F$ पूर्ण करणाऱ्या आणि $G$ पूर्ण करणाऱ्या वस्तूंची संख्या नेमकी समान आहे; म्हणजे $F$ पूर्ण करणाऱ्या व $G$ पूर्ण करणाऱ्या वस्तूंमध्ये द्विएक संबंध लावता येतो; म्हणजेच:
\begin{align*}
  \exists R(&\forall v\forall y(Rvy \lif (Fv \land Gy)) \land {}\\
    &\forall v(Fv \lif \lexists![y][Rvy]) \land {}\\
    &\forall y(Gy \lif \lexists![v][Rvy]))
\end{align*}
पण ह्यूमचे तत्त्व आणि कँटरचे तत्त्व यांत प्रकारभेद आहे. कँटरच्या तत्त्वातील चल “$A$” आणि “$B$” ही \emph{संच} दर्शवणारी प्रथम-क्रमाची पदे आहेत. ह्यूमच्या तत्त्वातील “$F$”, “$G$” आणि “$R$” ही प्रथम-क्रमाची पदे \emph{नाहीत}; ती \emph{विधेयस्थानी} येतात. (कदाचित ती \emph{गुणधर्म} दर्शवत असतील.) म्हणून ह्यूमचे तत्त्व साध्या भाषेत असे सांगता येईल: $F$ पूर्ण करणाऱ्या आणि $G$ पूर्ण करणाऱ्या वस्तूंची संख्या समान असते तेव्हा आणि केवळ तेव्हाच जेव्हा $F$ पूर्ण करणाऱ्या व $G$ पूर्ण करणाऱ्या वस्तूंमध्ये द्विएक संबंध असतो. याला \emph{ह्यूमचे तत्त्व} म्हणतात, कारण ह्यूम यांनी एकदा असे लिहिले:
\begin{quote}
  दोन संख्या अशा रीतीने जोडलेल्या असतील की एकीतील प्रत्येक एककाला दुसरीतील नेहमी एक एकक जुळते, तर आपण त्या समान आहेत असे म्हणतो.
  \citep[भाग III, पुस्तक 1, \S1]{Hume1740}
\end{quote}
फ्रेगे यांनी \citet[\S63]{Frege1884} मध्ये ह्यूम यांचा हा उतारा उद्धृत केला आणि नंतर (वस्तुतः) आपण ज्याला ह्यूमचे तत्त्व म्हटले आहे त्याची रूपरेषा मांडली; त्यामुळे आधुनिक गणितीय तर्कशास्त्रात या तत्त्वाला महत्त्व प्राप्त झाले.

ह्यूमचे तत्त्व आणि मूलनियम पाच यांतील रचनात्मक साम्य लक्षात घ्या. \olref[story][blv]{sec} मध्ये आपण मूलनियम पाच असे मांडले होते:
\[
  \fregeext{x}{F(x)} = \fregeext{x}{G(x)} \text{तेव्हा आणि केवळ तेव्हाच } \lforall[x][(F(x) \liff G(x))].
\]
या टप्प्यावर थोडे स्पष्टीकरण आणि तुलना उपयोगी ठरेल.

ह्यूमचे तत्त्व किंवा मूलनियम पाच यांसारख्या तत्त्वांचा अर्थ दोन प्रकारे घेता येतो: \emph{स्वसमावेशरहित} किंवा \emph{स्वसमावेशक} (\olref[story][predicative]{sec} आठवा). मूलनियम पाचच्या स्वसमावेशक अर्थानुसार, प्रत्येक~$F$ साठी $\fregeext{x}{F(x)}$ ही वस्तू, मूलनियम पाच मांडताना आपण ज्यावर संख्यापन केले त्याच प्रांतात येते. त्याचप्रमाणे ह्यूमच्या तत्त्वाच्या स्वसमावेशक अर्थानुसार, प्रत्येक~$F$ साठी $\fregenum{x}{F(x)}$ ही वस्तू, ह्यूमचे तत्त्व मांडताना ज्यावर संख्यापन केले त्याच प्रांतात येते. याउलट, \emph{स्वसमावेशरहित} अर्थानुसार $\fregeext{x}{F(x)}$ आणि~$\fregenum{x}{F(x)}$ या वस्तू एखाद्या \emph{वेगळ्या} प्रांतातील असतील.

मूलनियम पाचचा स्वसमावेशक अर्थ घेतल्यास सहज गुणधर्माधारित संचरचनेमुळे विसंगती निर्माण होते (तपशीलांसाठी \olref[story][blv]{sec} पाहा). सहज गुणधर्माधारित संचरचनेप्रमाणेच, त्याचा \emph{स्वसमावेशरहित} अर्थ घेतल्यास तो सुसंगत करता येतो. पण मग त्याच्याकडून अपेक्षित सर्व काही बहुधा साध्य होणार नाही.

ह्यूमचे तत्त्व मात्र सुसंगतपणे स्वसमावेशक अर्थाने घेता \emph{येते}. असा अर्थ घेतल्यावर ते बरेच सामर्थ्यवान ठरते.

उदाहरणार्थ, “$x \neq x$” हे विधेय पाहा; स्पष्टच कोणतीही वस्तू ते पूर्ण करत नाही. ह्यूमच्या तत्त्वावरून आता $\# x( x\neq
x)$ ही वस्तू मिळते. तिला आपण~$0$ ही संख्या मानू शकतो. आता \emph{स्वसमावेशक} अर्थानुसार—पण \emph{फक्त} त्याच अर्थानुसार—ही वस्तू $0$ आपल्या मूळ संख्यापन-प्रांतात येते. म्हणून “$x = 0$” या विधेयाला ह्यूमचे तत्त्व लावून $\#x (x =
0)$ ही वस्तू मिळवता येते. तिला आपण~$1$ ही संख्या मानू शकतो. शिवाय ह्यूमच्या तत्त्वावरून $0 \neq 1$ मिळते, कारण स्वतःशी असमान वस्तू आणि $0$ शी समान वस्तू यांच्यात द्विएक संबंध असू शकत नाही (पहिल्या प्रकारची एकही नाही, दुसऱ्या प्रकारची एक आहे). आता पुन्हा स्वसमावेशक अर्थाने $1$ आपल्या मूळ संख्यापन-प्रांतात येतो. म्हणून “$(x = 0 \lor x = 1)$” या विधेयाला ह्यूमचे तत्त्व लावून $\#x(x = 0 \lor
x = 1)$ ही वस्तू मिळते. तिला आपण~$2$ ही संख्या मानू शकतो; मग $0\neq 2$ आणि $1 \neq 2$ वगैरे दाखवता येते.

थोडक्यात, ह्यूमचे तत्त्व स्वसमावेशक अर्थाने घेतल्यास \emph{अनंत संख्येने वस्तू आहेत} हे त्यावरून मिळते. त्यामुळे \emph{नव-फ्रेगेवादी तर्कमीमांसावाद्यांना} अंकगणिताचा पाया म्हणून ह्यूमचे तत्त्व घ्यावेसे वाटले आहे.

पण फ्रेगे यांनी \emph{स्वतः} ह्यूमचे तत्त्व अंकगणिताचा पाया म्हणून घेतले नाही. उलट, एका स्पष्ट व्याख्येवरून त्यांनी ह्यूमचे तत्त्व सिद्ध केले: $\fregenum{x}{F(x)}$ ची व्याख्या $F \sim \Phi$ या संकल्पनेचा विस्तार अशी केली जाते. आधुनिक संज्ञेत हे $\fregenum{x}{F(x)} = \Setabs{G}{F \sim G}$ असे मांडण्याचा प्रयत्न करता येईल; पण मग सहज गुणधर्माधारित संचरचनेतील अडचणी पुन्हा उभ्या राहतील.

\end{document}
